\section{A novelty roadmap: what the work must become}
\label{sec:novelty-roadmap}

\subsection{Separate an engineering milestone from a research contribution}
A functioning simulator, an integrated ROS graph, a passing unit test or an
isolated delivery is an engineering milestone. It can be essential to research
without being a new scientific result. Likewise, correcting a stale timestamp
or a missing subscription is necessary experimental hygiene; it becomes a
research contribution only when embedded in a generalizable method or a
controlled finding about the effect of that mechanism.

The route to novelty should consequently be expressed as three questions:
\begin{enumerate}
\item \textbf{What does the nearest prior work already solve?} Name the closest
method and its assumptions, rather than comparing only with a weak default.
\item \textbf{What precisely changes?} Specify the new estimator, uncertainty
construction, decision rule, coordination contract or experimental intervention.
\item \textbf{What observation would demonstrate or refute its value?} Define
held-out outcomes and a fair alternative that could plausibly match the result.
\end{enumerate}
The following roadmap is proposed work. It does not report completed acceptance,
guarantee a publication, or certify novelty throughout literatures outside the
reviewed corpus.

\subsection{Primary contribution target: a joint observation--execution model}
\label{goal:primary}
\textbf{Intended contribution.} A measured, policy-conditioned short-horizon
motion envelope combined with acquisition-age-aware pedestrian occupancy, used
consistently to select and supervise commands on a walking humanoid.

The prior-art boundary is explicit. DWA already reasons about reachable and
stoppable velocities; the biped MPC paper already combines balance, uncertainty
and stopping; ASAP already learns a dynamics correction; and recent LiDAR policies
already perform sensor-driven avoidance \citep{dwa,biped2016,asap,safehumanoid,omni}.
The proposed method would need to contribute the \emph{particular measured
interface and decision construction}, not the generic combination of their keywords.

One candidate formulation represents executable robot occupancy as
\begin{equation}
\mathcal{R}_{\theta}(\tau\mid x,u,c)
=\widehat{\mathcal{R}}_{\theta}(\tau\mid x,u,c)
\oplus\mathcal{E}_{\theta}(\tau,c),
\end{equation}
where $c$ identifies policy/configuration and relevant posture/load conditions,
and $\mathcal{E}_{\theta}$ is a residual envelope evaluated on held-out response
trajectories. A person's predicted occupancy is aligned from the observation's
acquisition time, not reset by republishes. Candidate acceptance then concerns
overlap between the two occupied regions over the executable horizon. This is a
\emph{research specification}, not an already proved guarantee; valid coverage,
dependence assumptions and feasible fallback behavior must be established.

\textbf{Goal N1.1: characterize the interface.} Freeze the official or explicitly
selected payload policy, gains, filtering and command bounds. Measure steps,
ramps, stops, reversals, pivots and combined-axis transitions from varied initial
motion. Identify deadbands, asymmetric gains, delay, stopping drift and
cross-axis effects. Split calibration and test data by complete trajectories and
conditions rather than adjacent samples.

\textbf{Goal N1.2: define a parsimonious model.} Compare instantaneous dynamics,
generic bounded acceleration, the current hand-tuned envelope and a measured
model with and without calibrated residual uncertainty. Justify additional
parameters by held-out prediction or decision value. A more complex model that
only fits its training motions is not progress toward this contribution.

\textbf{Goal N1.3: demonstrate a downstream advantage.} Under matched sensing,
policy, actor schedule, goals and computational budget, show either lower contact
risk at similar completion or improved collision-free completion at comparable
risk. Report the frontier across conservatism settings instead of choosing one
favorable threshold. An equally conservative fixed-margin baseline is a required
alternative, because its success would refute the need for the extra model.

\subsection{Secondary contribution target: causal diagnosis of the onboard gap}
\label{goal:diagnosis}
\textbf{Intended contribution.} A controlled, source-audited decomposition of
how sensing, time alignment and execution cause navigation failures on the same
walking system. Benchmark platforms and perception datasets already exist
\citep{socnavbench,jrdb,metaurban,omninav}; the contribution must be new causal
findings or a validated protocol that makes previously confounded effects
measurable.

\textbf{Goal N2.1: define matched interventions.} Independently substitute
person observations, robot pose and map information while preserving evaluation
truth. Hold goals and actor exposure fixed when testing an algorithmic change;
otherwise a source mode may change which encounters occur. Add controlled age,
dropout, center bias, association errors and pose drift, with identities that
allow a consumed observation to be linked to a resulting action.

\textbf{Goal N2.2: measure interactions.} A small factorial design should test
response modeling against age correction and uncertainty treatment. Do not assume
that errors add independently. A result that tracking age dominates predictor
complexity, or that locomotion response dominates map error, is potentially useful
only if the experiment isolates that effect and survives held-out conditions.

\textbf{Goal N2.3: publish explanatory failures.} For representative contacts,
timeouts and invalid trials, reconstruct observation, estimate, proposal, veto,
policy-consumed command and measured response. Report the frequency of the
identified mechanism, not only one attractive anecdote. If source substitutions
change the accepted goals or crowd schedule, classify the experiment as
confounded until it is redesigned.

\subsection{Third contribution target: consistent supervision with liveness}
\label{goal:liveness}
\textbf{Intended contribution.} A decision or handoff rule that keeps route
following, social planning and final command supervision consistent about the
motion they are checking, and measurably reduces prolonged veto loops without
increasing contact or fall risk.

Freezing, cooperation and hierarchical navigation already have substantial prior
art \citep{freezing,rgl,biped2016,wheeledleg}. A new state machine with names such
as yield, wait and recover is not enough. The contribution would need a defined
invariant or optimization criterion, a bounded action-lifecycle contract, and
evidence explaining why the rule changes the tradeoff.

\textbf{Goal N3.1: compare against shared and independent shields.} Keep the
underlying predictor and response model fixed while varying coordination. If a
common final filter is attached to every planner, explain that the experiment
tests planners under a shared shield. If the shield is the independent variable,
do not let another hidden filter erase the difference.

\textbf{Goal N3.2: measure actual liveness.} Include repeated-veto duration,
command reversals, time to resume, useful progress after yielding, route age,
interventions and complete goals. A stationary robot can still be struck by an
approaching person; it is not a universally safe fallback. A controller that
avoids contact by remaining blocked forever also fails the intended task.

\textbf{Goal N3.3: test adverse scheduling.} Evaluate cancellation during
inference, late results, sensor starvation, competing command sources and changes
of mission generation. Some fixes here are baseline correctness rather than
novelty, but omitting them can invalidate the experiments used to support a new
coordination rule.

\subsection{Extension target: body, arm and payload risk}
\label{goal:body}
\textbf{Intended contribution.} Establish when a base-only footprint/envelope
is insufficient for a humanoid, and develop a sufficiently efficient approximation
of articulated swept occupancy if the evidence demonstrates a need.

Whole-body motion and capability-aware navigation already exist
\citep{expressive,asap,parkour,wheeledleg}. The first goal is therefore an
\emph{abstraction-validity study}: measure disagreement between the planner's
base approximation and link/payload clearance or native contacts. Include turns,
lateral motion, reach/carry transitions and recovery. Distinguish body contact,
object contact, fall and drop from personal-space intrusion.

If the simpler approximation is adequate in the declared conditions, the result
supports keeping it. If it misses important risks, compare a posture-conditioned
envelope with full articulated checking at matched computational budgets.
Demonstrating a useful fidelity--latency tradeoff would be a stronger contribution
than replacing every check with an expensive model without measuring its effect.
Magnetic fixture completion remains separate from physical finger grasping.

\subsection{Extension target: contextual human validation}
\label{goal:comfort}
\textbf{Intended contribution.} Determine whether a response-aware humanoid
motion strategy improves perceived predictability or comfort in a defined
interaction, beyond changes explained by speed and clearance.

SCAND already reports human evaluation, and the guidelines and proxemics taxonomy
provide substantial methodology \citep{scand,guidelines,proxemics2026}. The new
question must therefore be humanoid/task-specific, not ``add a questionnaire''.
Start with a bounded, counterbalanced study using matched encounter geometry.
Separate morphology, payload, yielding style, motion smoothness and clearance
where possible. Use participant-level analysis, declare the instrument and
population, and assess objective task outcomes alongside subjective responses.

Simulation video or human-in-the-loop avatars can be an intermediate instrument,
but their relationship to physical encounters must be stated
\citep{habitat3}. A gain in geometric clearance alone cannot close this research
goal. Equally, a favorable perception rating cannot replace contact and stability
evidence.

\subsection{A staged evidence ladder for the first paper}
\begin{longtable}{P{0.13\textwidth}P{0.36\textwidth}P{0.43\textwidth}}
\toprule Gate & Deliverable & Evidence required before the claim advances \\
\midrule\endfirsthead
\toprule Gate & Deliverable & Evidence required before the claim advances \\
\midrule\endhead
\bottomrule\endfoot
E0: honest baseline & Frozen source/policy/configuration and functioning scorer & Freshness, source and contact definitions verified; failed and invalid trials retained. \\
E1: measured model & Held-out command-response dataset & Prediction error and coverage by horizon, initial motion and operating condition; calibration kept separate from test. \\
E2: specific method & Reproducible model/decision specification & Equations, uncertainty assumptions, fallback behavior and timing budget; nearest methods identified. \\
E3: attribution & Matched baselines and ablations & Equal information and unavoidable execution limits; generic bounds and fixed margins can challenge the proposal. \\
E4: generalization & Multi-scenario and held-out evaluation & Complete-mission outcomes with denominators, uncertainty, repeated runs and evidence against scene-specific tuning. \\
E5: explanation & Failure taxonomy and trace examples & Mechanism frequency and representative counterexamples, not only selected successes. \\
E6: bounded claim & Manuscript and reproducibility package & A claim limited to demonstrated conditions, with unresolved access/literature gaps and negative results disclosed. \\
\end{longtable}

This ladder is a research-evidence sequence, not a replacement for the engineering
S0--S7 project plan. Code availability, a clean build and regression coverage
support E0. They do not automatically satisfy E1--E6. Likewise, an all-pass
development matrix is a useful milestone but does not by itself demonstrate
novelty or statistically strong generalization.

\subsection{What would make the first submission defensible?}
The strongest near-term submission should have one central contribution rather
than claim every gap at once. A sensible scope is N1 supported by the diagnostic
tools of N2, with N3 included only if a well-defined coordination mechanism is
actually necessary to realize the model's benefit. Full-body and human-evaluation
extensions can then form later studies with their own evidence requirements.

Before submission, extend the nearest-prior-work search specifically for
uncertainty-aware biped MPC, delayed/sampled-data safety filters, learned
reachable sets, decision-focused prediction and uncertainty calibration. Compare
the proposed mathematical object, not just method names. The four inaccessible
originals in this review should also be obtained where possible. A claim of
first-ever capability is unnecessary and fragile; a clearly demonstrated
conditional improvement or explanatory negative result is more defensible.

The intended result should ultimately be expressible in one sentence with a
declared comparator and scope: \emph{under matched sensing, controller and crowd
conditions, the proposed measured observation--execution model improves the
collision-free completion tradeoff, and controlled ablations identify when that
improvement depends on age, response calibration or uncertainty.} If the data do
not support that sentence, report which mechanism failed and revise the research
claim rather than loosening the evaluation definition.
